\documentclass[11pt]{article}

\usepackage[utf8]{inputenc}
\usepackage[T1]{fontenc}
\usepackage[spanish]{babel}
\usepackage{amsmath,amssymb,mathtools,bm,graphicx,xcolor}
\usepackage{float}
\usepackage{svg}
\svgsetup{inkscapelatex=false}
\usepackage[a4paper,margin=2.5cm]{geometry}
\usepackage[
    colorlinks=true,
    linkcolor=red,
    citecolor=green,
    urlcolor=red
]{hyperref}

\spanishdecimal{.}
\setlength{\parindent}{0pt}
\setlength{\parskip}{0.45em}
\setlength{\abovedisplayskip}{6pt}
\setlength{\belowdisplayskip}{6pt}
\setlength{\abovedisplayshortskip}{4pt}
\setlength{\belowdisplayshortskip}{4pt}

% Formato de encabezado y problemas
\newcommand{\encabezado}[3]{%
{\Large\bfseries #1\par}
\vspace{2pt}
{\normalsize #2\hfill #3\par}
\vspace{6pt}\hrule\vspace{8pt}}
\newcommand{\problema}[2]{%
\vspace{0.55em}
\noindent\textbf{#1.}\enspace{\small #2}\par
\vspace{0.3em}}
\newcommand{\inciso}[1]{%
\vspace{0.35em}
\noindent\textbf{(#1)}\enspace}

% Comandos auxiliares
\newcommand{\dd}{\,\mathrm{d}}
\newcommand{\ii}{\mathrm{i}}
\newcommand{\e}{\mathrm{e}}
\newcommand{\sen}{\operatorname{sen}}
\newcommand{\grad}{\nabla}
\newcommand{\laplaciano}{\nabla^2}
\newcommand{\R}{\mathbb{R}}

\renewcommand{\theequation}{7.\arabic{equation}}

\begin{document}

\encabezado{Funciones de Green: formas explícitas}{Física Matemática II}{J.R., G.R.}

\setcounter{equation}{29}

\problema{0}{Recordatorio: Laplaciano en coordenadas polares y esféricas}

Para una función escalar $F=F(\rho,\varphi)$, el operador Laplaciano en coordenadas polares se escribe como
\begin{equation*}
\nabla^2F=
\frac{1}{\rho}\frac{\partial\ }{\partial\rho}
\left(\rho\frac{\partial F}{\partial\rho}\right)
+\frac{1}{\rho^2}\frac{\partial^2F}{\partial\varphi^2}.
\end{equation*}
Si $F$ es radial, es decir $F=F(\rho)$, la parte angular se anula y queda
\begin{equation*}
\nabla^2F=
\frac{1}{\rho}\frac{\dd\ }{\dd\rho}
\left(\rho\frac{\dd F}{\dd\rho}\right).
\end{equation*}

Para una función escalar $F=F(r,\theta,\varphi)$, el operador Laplaciano en coordenadas esféricas se escribe como
\begin{equation*}
\nabla^2F=
\frac{1}{r^2}\frac{\partial\ }{\partial r}
\left(r^2\frac{\partial F}{\partial r}\right)
+\frac{1}{r^2\sen\theta}\frac{\partial\ }{\partial\theta}
\left(\sen\theta\frac{\partial F}{\partial\theta}\right)
+\frac{1}{r^2\sen^2\theta}\frac{\partial^2F}{\partial\varphi^2}.
\end{equation*}
Si $F$ es radial, es decir $F=F(r)$, la parte angular se anula y queda
\begin{equation*}
\nabla^2F=
\frac{1}{r^2}\frac{\dd\ }{\dd r}
\left(r^2\frac{\dd F}{\dd r}\right).
\end{equation*}

La función de Green definida por la ecuación de Green no es única. Si $G_1(\vec{x},\vec{x}')$ es solución, entonces $G_2(\vec{x},\vec{x}'):=G_1(\vec{x},\vec{x}')+H(\vec{x},\vec{x}')$ también es solución, si $H(\vec{x},\vec{x}')$ es solución del problema homogéneo asociado,
\begin{equation}
\hat{L}H(\vec{x},\vec{x}')=0.
\end{equation}
Este hecho permite considerar soluciones particulares simples como base para otras funciones de Green, que pueden obtenerse agregando una solución de la ecuación homogénea de modo que la función resultante satisfaga las condiciones de borde que simplifiquen el problema.

\problema{1}{Operador Laplaciano en $D=3$}

\textbf{Solución:} En este caso $p(\vec{x})=1$ y $q(\vec{x})=0$. En Física, usualmente se busca la función de Green que ``respete la homogeneidad e isotropía del espacio''. La primera condición implica que $G$ depende de $\vec{x}$ y $\vec{x}'$ sólo a través de su diferencia $\vec{x}-\vec{x}'$. La segunda condición implica que $G$ sólo depende del módulo de $\vec{x}-\vec{x}'$. Entonces
\begin{equation}
G(\vec{x},\vec{x}')=G(|\vec{x}-\vec{x}'|).
\end{equation}

En coordenadas esféricas centradas en $\vec{x}'$, escribimos $r:=|\vec{x}-\vec{x}'|$. Como $G=G(r)$, se tiene
\begin{equation}\label{n2Gd}
\nabla^2G=\frac{1}{r^2}\frac{\dd\ }{\dd r}\left(r^2\frac{\dd G}{\dd r}\right)=\delta^{(3)}(r).
\end{equation}
Por lo tanto, fuera del origen,
\begin{equation}\label{lap3_hom}
\frac{1}{r^2}\frac{\dd\ }{\dd r}\left(r^2\frac{\dd G}{\dd r}\right)=0,\qquad r\neq 0.
\end{equation}
La solución para $r\neq 0$ se encuentra integrando directamente \eqref{lap3_hom}. En efecto,
\begin{align*}
\frac{\dd\ }{\dd r}\left(r^2\frac{\dd G}{\dd r}\right)&=0,\\
r^2\frac{\dd G}{\dd r}&=C_1,\\
\frac{\dd G}{\dd r}&=\frac{C_1}{r^2}.
\end{align*}
Integrando una vez más,
\begin{equation*}
G(r)=C_2-\frac{C_1}{r}.
\end{equation*}
Renombrando constantes, obtenemos
\begin{equation}\label{lap3_general}
G(r)=\alpha+\frac{\beta}{r}.
\end{equation}

Por otro lado, la condición \eqref{n2Gd} implica, usando el teorema de Gauss, que
\begin{equation}\label{gauss_lap3}
\int_{\partial V}\vec\nabla G\cdot \dd\vec{S}
=\int_{\partial V}\frac{\dd G}{\dd r}r^2\dd\Omega=1.
\end{equation}
Desde \eqref{lap3_general},
\begin{equation*}
\frac{\dd G}{\dd r}=-\frac{\beta}{r^2}.
\end{equation*}
Luego, sobre una esfera de radio $R$,
\begin{align*}
\int_{\partial V}\frac{\dd G}{\dd r}r^2\dd\Omega
&=\int_{\partial V}\left(-\frac{\beta}{R^2}\right)R^2\dd\Omega\\
&=-\beta\int\dd\Omega\\
&=-4\pi\beta.
\end{align*}
La condición \eqref{gauss_lap3} requiere entonces que $-4\pi\beta=1$, o sea
\begin{equation*}
\beta=-\frac{1}{4\pi}.
\end{equation*}
Finalmente, en la mayoría de los problemas es conveniente elegir una función de Green que se anule para distancias muy grandes, es decir, tal que $\lim_{r\to\infty}G(r)=0$. Esta condición impone que $\alpha=0$ y por lo tanto
\begin{equation}
\boxed{G(\vec{x}-\vec{x}')=-\frac{1}{4\pi}\frac{1}{|\vec{x}-\vec{x}'|}.}
\end{equation}

\problema{2}{Operador Laplaciano en $D=2$}

\textbf{Solución:} En coordenadas polares $(\rho,\varphi)$ centradas en $\vec{x}'$, escribimos $\rho:=|\vec{x}-\vec{x}'|$. Como $G=G(\rho)$,
\begin{equation}\label{lap2_delta}
\nabla^2G=\frac{1}{\rho}\frac{\dd\ }{\dd\rho}\left(\rho\frac{\dd G}{\dd\rho}\right)=\delta^{(2)}(\rho).
\end{equation}
Integrando la ecuación para $\rho\neq0$, es decir,
\begin{equation}\label{lap2_hom}
\frac{1}{\rho}\frac{\dd\ }{\dd\rho}\left(\rho\frac{\dd G}{\dd\rho}\right)=0,
\end{equation}
encontramos
\begin{align*}
\frac{\dd\ }{\dd\rho}\left(\rho\frac{\dd G}{\dd\rho}\right)&=0,\\
\rho\frac{\dd G}{\dd\rho}&=C_1,\\
\frac{\dd G}{\dd\rho}&=\frac{C_1}{\rho},\\
G(\rho)&=C_2+C_1\ln\rho.
\end{align*}
Renombrando constantes,
\begin{equation}\label{lap2_general}
G(\rho)=\alpha+\beta\ln\rho.
\end{equation}

Nuevamente, el teorema de Gauss, en su versión 2D, implica
\begin{equation}\label{gauss_lap2}
\int_{\partial S}\vec\nabla G\cdot \dd\vec{S}
=\oint\frac{\dd G}{\dd\rho}\rho\,\dd\varphi=2\pi\beta=1.
\end{equation}
La última igualdad se obtiene directamente desde \eqref{lap2_general}, pues
\begin{equation*}
\frac{\dd G}{\dd\rho}=\frac{\beta}{\rho},
\end{equation*}
y por lo tanto
\begin{align*}
\oint\frac{\dd G}{\dd\rho}\rho\,\dd\varphi
&=\oint\frac{\beta}{\rho}\rho\,\dd\varphi\\
&=\beta\int_0^{2\pi}\dd\varphi\\
&=2\pi\beta.
\end{align*}
Así, $2\pi\beta=1$, de donde $\beta=1/(2\pi)$. De este modo, eliminando el término constante, que es una solución de la ecuación homogénea, encontramos
\begin{equation}
\boxed{G(\vec{x}-\vec{x}')=\frac{1}{2\pi}\ln{|\vec{x}-\vec{x}'|}.}
\end{equation}

\problema{3}{Operador de Helmholtz en $D=3$}

\textbf{Solución:} En el caso del operador de Helmholtz $\hat{L}=\nabla^2+k^2$, corresponde tomar $p(\vec{x})=1$ y $q(\vec{x})=k^2$. Buscaremos la función de Green de la forma
\begin{equation*}
G(\vec{x},\vec{x}')=G(|\vec{x}-\vec{x}'|).
\end{equation*}

En coordenadas esféricas centradas en $\vec{x}'$, con $r:=|\vec{x}-\vec{x}'|$,
\begin{equation}\label{HGd}
(\nabla^2+k^2)G
=\frac{1}{r^2}\frac{\dd\ }{\dd r}\left(r^2\frac{\dd G}{\dd r}\right)+k^2G
=\delta^{(3)}(r),
\end{equation}
por lo tanto, para $r\neq0$,
\begin{equation}\label{HG0}
\frac{1}{r^2}\frac{\dd\ }{\dd r}\left(r^2\frac{\dd G}{\dd r}\right)+k^2G=0,\qquad r\neq0.
\end{equation}
Para solucionar \eqref{HG0} realizamos el cambio de variable
\begin{equation*}
G(r)=\frac{u(r)}{r}.
\end{equation*}
Entonces
\begin{align*}
\frac{\dd G}{\dd r}
&=\frac{\dd\ }{\dd r}\left(\frac{u}{r}\right)\\
&=\frac{u'}{r}-\frac{u}{r^2}\\
&=\frac{ru'-u}{r^2},
\end{align*}
y por lo tanto
\begin{equation*}
r^2\frac{\dd G}{\dd r}=ru'-u.
\end{equation*}
Derivando,
\begin{align*}
\frac{\dd\ }{\dd r}\left(r^2\frac{\dd G}{\dd r}\right)
&=\frac{\dd\ }{\dd r}(ru'-u)\\
&=u'+ru''-u'\\
&=ru''.
\end{align*}
Con esto, \eqref{HG0} se reduce a
\begin{equation*}
\frac{1}{r^2}(ru'')+k^2\frac{u}{r}=0,
\end{equation*}
es decir,
\begin{equation*}
u''+k^2u=0.
\end{equation*}
Por lo tanto, las soluciones son de la forma
\begin{equation}\label{GH3d}
G(r)=\alpha\frac{\e^{\ii kr}}{r}+\beta\frac{\e^{-\ii kr}}{r}.
\end{equation}

Con esta solución, válida para $r\neq0$, retornamos a la ecuación \eqref{HGd} que, luego de integrar sobre una esfera de radio $R$ centrada en $r=0$ y usando el teorema de Gauss, implica que
\begin{align}
1 &= \oint_{\partial V}\vec\nabla G\cdot \dd\vec{S}+k^2\int_VG\,\dd V \label{helm_int1}\\
&= \oint_{\partial V}\frac{\dd G}{\dd r}r^2\dd\Omega+k^2\int_VGr^2\,\dd r\dd\Omega \label{helm_int2}\\
&= 4\pi G'R^2+4\pi k^2\int_0^R Gr^2\,\dd r \label{helm_int3}\\
&= 4\pi\left[\alpha\left(\ii kR-1\right)\e^{\ii kR}-\beta\left(\ii kR+1\right)\e^{-\ii kR}\right]
+4\pi k^2\int_0^R \left[\alpha\e^{\ii kr}+\beta\e^{-\ii kr}\right]r\,\dd r \label{helm_int4}\\
&= -4\pi(\alpha+\beta). \label{helm_int5}
\end{align}
En \eqref{helm_int4} se usó que, desde \eqref{GH3d},
\begin{align*}
\frac{\dd G}{\dd r}
&=\alpha\frac{\dd\ }{\dd r}\left(\frac{\e^{\ii kr}}{r}\right)
+\beta\frac{\dd\ }{\dd r}\left(\frac{\e^{-\ii kr}}{r}\right)\\
&=\alpha\frac{(\ii kr-1)\e^{\ii kr}}{r^2}
-\beta\frac{(\ii kr+1)\e^{-\ii kr}}{r^2}.
\end{align*}
Además,
\begin{align*}
k^2\int_0^R \e^{\ii kr}r\,\dd r &= \e^{\ii kR}(1-\ii kR)-1,\\
k^2\int_0^R \e^{-\ii kr}r\,\dd r &= \e^{-\ii kR}(1+\ii kR)-1.
\end{align*}
Por lo tanto, el término de borde más el término de volumen dentro de \eqref{helm_int4} queda
\begin{align*}
&\alpha(\ii kR-1)\e^{\ii kR}-\beta(\ii kR+1)\e^{-\ii kR}\\
&\qquad +\alpha\left[\e^{\ii kR}(1-\ii kR)-1\right]
+\beta\left[\e^{-\ii kR}(1+\ii kR)-1\right]\\
&=-(\alpha+\beta),
\end{align*}
lo cual da el paso final de \eqref{helm_int4} a \eqref{helm_int5}.

De esta forma encontramos la condición
\begin{equation}\label{cond_ab}
\alpha+\beta=-\frac{1}{4\pi},
\end{equation}
que permite escribir la solución \eqref{GH3d} como
\begin{equation}\label{GH3d2}
G(r)=-\frac{1}{4\pi}\frac{\e^{\ii kr}}{r}-2\ii\beta\frac{\sen(kr)}{r}.
\end{equation}
En efecto, usando \eqref{cond_ab}, se tiene $\alpha=-1/(4\pi)-\beta$, y entonces
\begin{align*}
G(r)
&=\left(-\frac{1}{4\pi}-\beta\right)\frac{\e^{\ii kr}}{r}+\beta\frac{\e^{-\ii kr}}{r}\\
&=-\frac{1}{4\pi}\frac{\e^{\ii kr}}{r}+\beta\frac{\e^{-\ii kr}-\e^{\ii kr}}{r}\\
&=-\frac{1}{4\pi}\frac{\e^{\ii kr}}{r}-2\ii\beta\frac{\sen(kr)}{r}.
\end{align*}
El segundo término es una solución de la ecuación de Helmholtz \textit{homogénea}; es proporcional a la función esférica de Bessel $j_0(kr)$. Por lo tanto, es posible elegir
\begin{equation}\label{GH3d3}
\boxed{G^{+}(\vec{x},\vec{x}')=-\frac{1}{4\pi}\frac{\e^{\ii k|\vec{x}-\vec{x}'|}}{|\vec{x}-\vec{x}'|},}
\end{equation}
que corresponde a la elección $\beta=0$. Análogamente, podemos considerar la función de Green
\begin{equation}\label{GH3d4}
G^{-}(\vec{x},\vec{x}')=-\frac{1}{4\pi}\frac{\e^{-\ii k|\vec{x}-\vec{x}'|}}{|\vec{x}-\vec{x}'|},
\end{equation}
que se encuentra en el caso en que $\beta=-1/4\pi$.

\end{document}
